\documentclass[11pt]{article}
\usepackage{preamble}

\newcommand{\lessonrow}[2]{\textbf{#1} & #2 \\ \hline}
\newcommand{\IconScope}{[SCOPE]}

\begin{document}

\packetheading{Lesson 4-3 · Period 1 · Student Packet}{Rational Expressions --- Equivalent \& Simplify}

\sectionbanner{OBJECTIVES}
{\small
\renewcommand{\arraystretch}{1.25}
\noindent\begin{tabular}{|p{1.8in}|p{4.8in}|}
\hline
\lessonrow{Math Objective}{Write equivalent rational expressions over a restricted domain, identify domain restrictions from factored denominators, and simplify rational expressions by canceling common factors.}
\lessonrow{Language Objective}{Explain, in writing and orally using sentence frames, why a value is excluded from the domain and why we cancel factors but not terms.}
\lessonrow{Essential Understanding}{A rational expression equals an equivalent simpler form ONLY when both share the same domain --- canceling a factor never erases the restriction it imposed.}
\end{tabular}
}

\sectionbanner{TOPIC GOALS / ESSENTIAL QUESTION / MATERIALS}
{\small
\renewcommand{\arraystretch}{1.25}
\noindent\begin{tabular}{|p{1.8in}|p{4.8in}|}
\hline
\lessonrow{Topic Goals}{Write equivalent rational expressions, simplify, and identify domain restrictions from polynomial denominators.}
\lessonrow{Essential Question}{Why does canceling a common factor not remove its domain restriction?}
\lessonrow{Materials}{Student packet, pencil, highlighter. Teacher displays worked Example 1 then gradual-release into Example 2.}
\end{tabular}
}

\sectionbanner{LAUNCH --- Rational Expressions (teacher models on screen)}

\begin{callout}{\IconBook\ EQUIVALENT + SIMPLIFY RULES (keep printed on your page)}{calloutgreen}
1. Factor the numerator AND denominator completely.

2. Cancel only common FACTORS (never terms across + or $-$).

3. State domain restrictions: every zero of the ORIGINAL denominator is excluded --- including factors you canceled.
\end{callout}

\begin{callout}{\IconWarn\ COMMON ERROR}{calloutred}
Do NOT cancel terms that are added or subtracted across a fraction bar.

You can only cancel factors that multiply across the entire numerator or denominator.

Test: factor completely FIRST, then cancel only identical factor pairs.
\end{callout}

\sectionbanner{EXPLORE --- Think-Pair-Share (33 min)}
\textit{Work with a partner. Use the Equivalent + Simplify Rules above for every problem. Move through items in order --- each builds on the last.}

\textbf{Bank-item labels (in order):}
\begin{enumerate}[leftmargin=1.6em,itemsep=2pt,topsep=2pt]
  \item Try It 1 --- Write an equivalent rational expression
  \item Practice \#20 --- Identify equivalent expressions
  \item Practice \#22 --- Identify domain restrictions
  \item Practice \#24 --- Apply domain restrictions
  \item Try It 2 --- Simplify by factoring and canceling
  \item Practice \#17 --- Simplify
  \item Practice \#19 --- Simplify with domain restrictions (extension)
\end{enumerate}

\bankitem{Try It 1 --- Write an equivalent rational expression}{%
Write an expression equivalent to \((x-4)/(x)\) over the domain \(\{x\mid x\ne 0 \text{ or } -2\}\).
}{}

\bankitem{Practice \#20 --- Identify equivalent expressions}{%
Write an equivalent expression over the given domain.

\((x(x-2))/(x-5)\) for all \(x\) except \(5\) and \(-6\)
}{}

\bankitem{Practice \#22 --- Identify domain restrictions}{%
What is the simplified form of the rational expression? What is the domain?

\((y^2-5y-24)/(y^2+3y)\)
}{}

\bankitem{Practice \#24 --- Apply domain restrictions}{%
What is the simplified form of the rational expression? What is the domain?

\((x^2+8x+15)/(x^2-x-12)\)
}{}

\bankitem{Try It 2 --- Simplify by factoring and canceling}{%
Simplify each expression and state the domain.

\begin{enumerate}[label=\textbf{\alph*.},leftmargin=1.8em,itemsep=3pt,topsep=4pt]
  \item \((x^2+2x+1)/(x^3-2x^2-3x)\)
  \item \((x^3+4x^2-x-4)/(x^2+3x-4)\)
\end{enumerate}
}{}

\bankitem{Practice \#17 --- Simplify}{%
Construct Arguments Explain how you can tell whether a rational expression is in simplest form.
}{}

\bankitem{Practice \#19 --- Simplify with domain restrictions (extension)}{%
Reason If the denominator of a rational expression is \(x^3+3x^2-10x\), what value(s) must be restricted from the domain for \(x\)?
}{}

\sectionbanner{SHARE / SUMMARY (5 min)}

\begin{sentenceframebox}
\textbf{Sentence frames}

Pair share: ``The factored form is \_\_\_, so the domain excludes \(x =\) \_\_\_. After canceling the factor \_\_\_, the simplified expression is \_\_\_ with the same domain restriction.''

Whole class: one pair reads their sentence. Class signals thumbs up/sideways.
\end{sentenceframebox}

\begin{summaryexitbox}{EXIT TICKET --- Summary Recap}
To simplify \((x^2+3x-10)/(x^2-4)\) I first factor to \_\_\_.

The restricted domain values are \(x \ne\) \_\_\_.

After canceling the factor \_\_\_, my simplified expression is \_\_\_.

One biggest thing I learned today: \_\_\_\_\_\_\_\_\_\_\_\_\_\_\_\_\_\_\_\_\_\_\_
\end{summaryexitbox}

\clearpage

\sectionbanner{OPTIONAL REINFORCEMENT (not collected)}
\textit{If you finish Explore early or want extra practice before Period 2.}

\begin{callout}{\IconScope\ OPTIONAL · Model \& Discuss}{calloutblue}
Return to the Do Now graph. With your partner:

B) Use Structure. Why can the graph never have a point at the skipped x-value, no matter what simplified form we write?

C) Examine at least two rational expressions you simplified in Explore. For each, explain in your own words why the domain restriction from the canceled factor still applies.
\end{callout}

\bankitem{Optional \#1  (Savvas Practice)}{%
What is the simplified form of the rational expression? What is the domain?

\((x^3+9x^2-10x)/(x^3-9x^2-10x)\)
}{}

\end{document}
